% Generated draft: requires mathematical and source review.
\documentclass[11pt]{article}
\usepackage[margin=25mm]{geometry}
\usepackage{fontspec}
\setmainfont{Latin Modern Roman}
\usepackage{amsmath,amssymb,mathtools,unicode-math}
\setmathfont{Latin Modern Math}
\usepackage{enumitem,needspace,xurl,hyperref}
\setlength{\parindent}{0pt}
\setlength{\parskip}{5pt}
\setlength{\emergencystretch}{3em}
\newcommand{\N}{\mathbb{N}}
\newcommand{\Z}{\mathbb{Z}}
\newcommand{\Q}{\mathbb{Q}}
\newcommand{\R}{\mathbb{R}}
\newcommand{\C}{\mathbb{C}}
\newcommand{\F}{\mathbb{F}}
\newcommand{\A}{\mathbb{A}}
\newcommand{\PP}{\mathbb{P}}
\newcommand{\E}{\mathbb{E}}
\newcommand{\eps}{\varepsilon}
\newcommand{\dd}{\,\mathrm{d}}
\begin{document}
\section*{UnsolvedMath ID: 10000078}
\textbf{Source problem code:} catalogue entry
\subsection{Definitions and mathematical statement}
For definiteness of notation, let \(V_n=\{1,\ldots,n\}^2\), \(n\ge2\), with nearest-neighbor edges. Write \(a_n=(1,n)\), \(b_n=(n,n)\), \(c_n=(1,1)\), and \(d_n=(n,1)\) for its top and bottom corners. Given a real field \(h_n\) on \(V_n\), assign vertex weight \(w_n(v)=\exp(h_n(v))\). With the usual vertex-sum convention, a path \(P=(v_0,\ldots,v_m)\) has length \(L_n(P)=\sum_{j=0}^{m}w_n(v_j)\), and
\[
 D_n(x,y)=\min_{P:x\to y}L_n(P)\qquad(x\ne y),\qquad D_n(x,x)=0.
\]
Choose shortest paths \(\gamma_1(n)\) from \(a_n\) to \(b_n\) and \(\gamma_2(n)\) from \(c_n\) to \(d_n\), with a specified tie rule if necessary. The first requested assertion is
\[
 \exists c>0\ \forall n\ge2:\quad
 \PP\bigl(V(\gamma_1(n))\cap V(\gamma_2(n))\ne\varnothing\bigr)>c.
\]
Here the field is described in the source as a Gaussian free field ``with no boundary conditions.'' On a finite connected graph the quadratic energy \(\sum_{\{u,v\}\in E}(h(u)-h(v))^2\) vanishes on constant fields, so a probability law with this free-boundary energy requires a choice such as a mean-zero representative, pinning, or a quotient by additive constants, as well as an energy normalization. Adding a common constant multiplies all weights and distances by the same positive factor and preserves the minimizing paths, but the record still does not specify the field law or normalized metric.

The remaining requests are to construct metric and geodesic scaling limits, determine geodesic dimension and height concentration, and identify any SLE parameter of limiting paths. The spatial and distance rescalings, convergence topology, definition of height, dimension convention, curve parametrization and boundary data are not supplied. These requests are retained as an unresolved program rather than assigned a particular Liouville metric or SLE law.
\subsection{Short English statement}
For the exponential of a Gaussian free field on a square grid, prove a positive lower bound, independent of grid size, for intersection of the top-to-top and bottom-to-bottom geodesics. The record also asks for scaling limits and their geometric properties. Its supplied partial-progress assessment discusses related GFF-weighted metrics without specifying enough data to identify these limits or an SLE parameter.
\subsection{Sources}
\begin{itemize}
\item Local UnsolvedMath registry, permanent ID \(10000078\), exact source code \texttt{catalogue entry}, ``Geodesics in Gaussian-free-field random metrics''; attributed to Itai Benjamini, no proposed year supplied.
\item Benjamini, \emph{Random Planar Metrics}, Section 3, item (5), PDF page 3: \url{https://www.wisdom.weizmann.ac.il/~itai/randomplanar2.pdf#page=3}. Preserved at \url{https://arquivo.pt/noFrame/replay/20201231041544id_/http://www.wisdom.weizmann.ac.il/~itai/randomplanar2.pdf}.
\item Supplied verification dated 2026-07-29: pypdf extraction checked against rendered PDF; retained after cross-document deduplication and source-side status review. Source evidence NEEDS\_REVIEW records an open or conjectural item and no complete primary-literature resolution in a bounded modern audit.
\item Supplied research classification PARTIAL-PROGRESS and registry status partially\_solved. The summary names Ding--Goswami, a ``negative moments''/``Liouville metric'' line, and ``GFF metric''/``LSAT with GFF weights,'' reporting work on scaling limits but unresolved crossing/geodesic intersections and SLE parameters. It supplies no complete bibliographic references for those developments.
\item Dataset attribution: Ulam AI / UnsolvedMath Contributors, CC BY 4.0. Source rights NEEDS\_REVIEW: publicly accessible author-hosted or publisher copy, with release reuse terms requiring review.
\end{itemize}
\end{document}
